% Figure: rows of the strain matrix E near a level transition and a no-slip
% wall (Chapter 9). Cells, faces, the rows of E that contain one fine
% segment and the complete shear row of a hanging vertex are read from the
% code's VariableStaggeredDiffusionOperator: generated/viscosity.tex
% (viscosity.jl, viscosity_data). Labels hand-written.
\begin{figure}
  \centering
  \def\fxwnone{%
    \def\fxwcell##1##2##3##4##5{}\def\fxwface##1##2##3##4##5{}%
    \def\fxwrow##1##2##3##4{}\def\fxwstencil##1##2##3##4{}\def\fxwhanging##1##2{}}%
  % Background of both panels: level shading, cells, wall, faces (red:
  % normal e_1, blue: e_2; the target segment of panel (a) thick).
  \def\fxwbackground#1{%
    \begin{scope}
      \clip (3.0,0.0) rectangle (5.0,1.25);
      \fxwnone
      \def\fxwcell##1##2##3##4##5{\fill[tr level ##5] (##1,##2) rectangle (##3,##4);}%
      % Generated by docs/techref/figures/generate.jl; do not edit.
\fxwcell{3}{0}{3.5}{0.5}{0}
\fxwcell{3.5}{0}{4}{0.5}{0}
\fxwcell{3}{0.5}{3.5}{1}{0}
\fxwcell{3.5}{0.5}{4}{1}{0}
\fxwcell{3}{1}{3.5}{1.5}{0}
\fxwcell{3.5}{1}{4}{1.5}{0}
\fxwcell{4}{0}{4.25}{0.25}{1}
\fxwcell{4.25}{0}{4.5}{0.25}{1}
\fxwcell{4.5}{0}{4.75}{0.25}{1}
\fxwcell{4.75}{0}{5}{0.25}{1}
\fxwcell{4}{0.25}{4.25}{0.5}{1}
\fxwcell{4.25}{0.25}{4.5}{0.5}{1}
\fxwcell{4.5}{0.25}{4.75}{0.5}{1}
\fxwcell{4.75}{0.25}{5}{0.5}{1}
\fxwcell{4}{0.5}{4.25}{0.75}{1}
\fxwcell{4.25}{0.5}{4.5}{0.75}{1}
\fxwcell{4.5}{0.5}{4.75}{0.75}{1}
\fxwcell{4.75}{0.5}{5}{0.75}{1}
\fxwcell{4}{0.75}{4.25}{1}{1}
\fxwcell{4.25}{0.75}{4.5}{1}{1}
\fxwcell{4.5}{0.75}{4.75}{1}{1}
\fxwcell{4.75}{0.75}{5}{1}{1}
\fxwcell{4}{1}{4.25}{1.25}{1}
\fxwcell{4.25}{1}{4.5}{1.25}{1}
\fxwcell{4.5}{1}{4.75}{1.25}{1}
\fxwcell{4.75}{1}{5}{1.25}{1}
\fxwface{1}{3}{0.25}{0.5}{0}
\fxwface{1}{3}{0.75}{0.5}{0}
\fxwface{1}{3}{1.25}{0.5}{0}
\fxwface{1}{3.5}{0.25}{0.5}{0}
\fxwface{1}{3.5}{0.75}{0.5}{0}
\fxwface{1}{3.5}{1.25}{0.5}{0}
\fxwface{1}{4}{0.125}{0.25}{1}
\fxwface{1}{4}{0.375}{0.25}{0}
\fxwface{1}{4}{0.625}{0.25}{0}
\fxwface{1}{4}{0.875}{0.25}{0}
\fxwface{1}{4}{1.125}{0.25}{0}
\fxwface{1}{4.25}{0.125}{0.25}{0}
\fxwface{1}{4.25}{0.375}{0.25}{0}
\fxwface{1}{4.25}{0.625}{0.25}{0}
\fxwface{1}{4.25}{0.875}{0.25}{0}
\fxwface{1}{4.25}{1.125}{0.25}{0}
\fxwface{1}{4.5}{0.125}{0.25}{0}
\fxwface{1}{4.5}{0.375}{0.25}{0}
\fxwface{1}{4.5}{0.625}{0.25}{0}
\fxwface{1}{4.5}{0.875}{0.25}{0}
\fxwface{1}{4.5}{1.125}{0.25}{0}
\fxwface{1}{4.75}{0.125}{0.25}{0}
\fxwface{1}{4.75}{0.375}{0.25}{0}
\fxwface{1}{4.75}{0.625}{0.25}{0}
\fxwface{1}{4.75}{0.875}{0.25}{0}
\fxwface{1}{4.75}{1.125}{0.25}{0}
\fxwface{1}{5}{0.125}{0.25}{0}
\fxwface{1}{5}{0.375}{0.25}{0}
\fxwface{1}{5}{0.625}{0.25}{0}
\fxwface{1}{5}{0.875}{0.25}{0}
\fxwface{1}{5}{1.125}{0.25}{0}
\fxwface{2}{4.125}{0.25}{0.25}{0}
\fxwface{2}{4.375}{0.25}{0.25}{0}
\fxwface{2}{4.625}{0.25}{0.25}{0}
\fxwface{2}{4.875}{0.25}{0.25}{0}
\fxwface{2}{3.25}{0.5}{0.5}{0}
\fxwface{2}{3.75}{0.5}{0.5}{0}
\fxwface{2}{4.125}{0.5}{0.25}{0}
\fxwface{2}{4.375}{0.5}{0.25}{0}
\fxwface{2}{4.625}{0.5}{0.25}{0}
\fxwface{2}{4.875}{0.5}{0.25}{0}
\fxwface{2}{4.125}{0.75}{0.25}{0}
\fxwface{2}{4.375}{0.75}{0.25}{0}
\fxwface{2}{4.625}{0.75}{0.25}{0}
\fxwface{2}{4.875}{0.75}{0.25}{0}
\fxwface{2}{3.25}{1}{0.5}{0}
\fxwface{2}{3.75}{1}{0.5}{0}
\fxwface{2}{4.125}{1}{0.25}{0}
\fxwface{2}{4.375}{1}{0.25}{0}
\fxwface{2}{4.625}{1}{0.25}{0}
\fxwface{2}{4.875}{1}{0.25}{0}
\fxwface{2}{4.125}{1.25}{0.25}{0}
\fxwface{2}{4.375}{1.25}{0.25}{0}
\fxwface{2}{4.625}{1.25}{0.25}{0}
\fxwface{2}{4.875}{1.25}{0.25}{0}
\fxwrow{0}{3.75}{0.25}{1}
\fxwrow{0}{4.125}{0.125}{-4}
\fxwrow{1}{4}{0}{8}
\fxwrow{1}{4}{0.25}{-4}
\fxwhanging{4}{0.25}
\fxwstencil{1}{4}{0.125}{-4}
\fxwstencil{1}{4}{0.375}{4}
\fxwstencil{2}{4.125}{0.25}{0.477}
\fxwstencil{2}{4.375}{0.25}{0.243}
\fxwstencil{2}{4.625}{0.25}{-0.0063}
\fxwstencil{2}{3.75}{0.5}{-1.08}
\fxwstencil{2}{4.125}{0.5}{0.235}
\fxwstencil{2}{4.375}{0.5}{0.143}
\fxwstencil{2}{4.625}{0.5}{-0.0107}
\fxwstencil{2}{4.125}{0.75}{0.144}
\fxwstencil{2}{4.375}{0.75}{0.094}
%
      \fxwnone
      \def\fxwcell##1##2##3##4##5{\draw[tr cell] (##1,##2) rectangle (##3,##4);}%
      \def\fxwface##1##2##3##4##5{%
        \ifnum##5=1 \ifnum#1=1
          \draw[line width=2.4pt,trred] (##2,##3-##4/2+0.012) -- (##2,##3+##4/2-0.012);\fi\fi
        \ifnum##1=1 \draw[trred!80!black,line width=0.5pt,-{Stealth[length=2.4pt,width=2pt]}] (##2-0.035,##3) -- (##2+0.035,##3);
        \else \draw[trblue!80!black,line width=0.5pt,-{Stealth[length=2.4pt,width=2pt]}] (##2,##3-0.035) -- (##2,##3+0.035);\fi}%
      % Generated by docs/techref/figures/generate.jl; do not edit.
\fxwcell{3}{0}{3.5}{0.5}{0}
\fxwcell{3.5}{0}{4}{0.5}{0}
\fxwcell{3}{0.5}{3.5}{1}{0}
\fxwcell{3.5}{0.5}{4}{1}{0}
\fxwcell{3}{1}{3.5}{1.5}{0}
\fxwcell{3.5}{1}{4}{1.5}{0}
\fxwcell{4}{0}{4.25}{0.25}{1}
\fxwcell{4.25}{0}{4.5}{0.25}{1}
\fxwcell{4.5}{0}{4.75}{0.25}{1}
\fxwcell{4.75}{0}{5}{0.25}{1}
\fxwcell{4}{0.25}{4.25}{0.5}{1}
\fxwcell{4.25}{0.25}{4.5}{0.5}{1}
\fxwcell{4.5}{0.25}{4.75}{0.5}{1}
\fxwcell{4.75}{0.25}{5}{0.5}{1}
\fxwcell{4}{0.5}{4.25}{0.75}{1}
\fxwcell{4.25}{0.5}{4.5}{0.75}{1}
\fxwcell{4.5}{0.5}{4.75}{0.75}{1}
\fxwcell{4.75}{0.5}{5}{0.75}{1}
\fxwcell{4}{0.75}{4.25}{1}{1}
\fxwcell{4.25}{0.75}{4.5}{1}{1}
\fxwcell{4.5}{0.75}{4.75}{1}{1}
\fxwcell{4.75}{0.75}{5}{1}{1}
\fxwcell{4}{1}{4.25}{1.25}{1}
\fxwcell{4.25}{1}{4.5}{1.25}{1}
\fxwcell{4.5}{1}{4.75}{1.25}{1}
\fxwcell{4.75}{1}{5}{1.25}{1}
\fxwface{1}{3}{0.25}{0.5}{0}
\fxwface{1}{3}{0.75}{0.5}{0}
\fxwface{1}{3}{1.25}{0.5}{0}
\fxwface{1}{3.5}{0.25}{0.5}{0}
\fxwface{1}{3.5}{0.75}{0.5}{0}
\fxwface{1}{3.5}{1.25}{0.5}{0}
\fxwface{1}{4}{0.125}{0.25}{1}
\fxwface{1}{4}{0.375}{0.25}{0}
\fxwface{1}{4}{0.625}{0.25}{0}
\fxwface{1}{4}{0.875}{0.25}{0}
\fxwface{1}{4}{1.125}{0.25}{0}
\fxwface{1}{4.25}{0.125}{0.25}{0}
\fxwface{1}{4.25}{0.375}{0.25}{0}
\fxwface{1}{4.25}{0.625}{0.25}{0}
\fxwface{1}{4.25}{0.875}{0.25}{0}
\fxwface{1}{4.25}{1.125}{0.25}{0}
\fxwface{1}{4.5}{0.125}{0.25}{0}
\fxwface{1}{4.5}{0.375}{0.25}{0}
\fxwface{1}{4.5}{0.625}{0.25}{0}
\fxwface{1}{4.5}{0.875}{0.25}{0}
\fxwface{1}{4.5}{1.125}{0.25}{0}
\fxwface{1}{4.75}{0.125}{0.25}{0}
\fxwface{1}{4.75}{0.375}{0.25}{0}
\fxwface{1}{4.75}{0.625}{0.25}{0}
\fxwface{1}{4.75}{0.875}{0.25}{0}
\fxwface{1}{4.75}{1.125}{0.25}{0}
\fxwface{1}{5}{0.125}{0.25}{0}
\fxwface{1}{5}{0.375}{0.25}{0}
\fxwface{1}{5}{0.625}{0.25}{0}
\fxwface{1}{5}{0.875}{0.25}{0}
\fxwface{1}{5}{1.125}{0.25}{0}
\fxwface{2}{4.125}{0.25}{0.25}{0}
\fxwface{2}{4.375}{0.25}{0.25}{0}
\fxwface{2}{4.625}{0.25}{0.25}{0}
\fxwface{2}{4.875}{0.25}{0.25}{0}
\fxwface{2}{3.25}{0.5}{0.5}{0}
\fxwface{2}{3.75}{0.5}{0.5}{0}
\fxwface{2}{4.125}{0.5}{0.25}{0}
\fxwface{2}{4.375}{0.5}{0.25}{0}
\fxwface{2}{4.625}{0.5}{0.25}{0}
\fxwface{2}{4.875}{0.5}{0.25}{0}
\fxwface{2}{4.125}{0.75}{0.25}{0}
\fxwface{2}{4.375}{0.75}{0.25}{0}
\fxwface{2}{4.625}{0.75}{0.25}{0}
\fxwface{2}{4.875}{0.75}{0.25}{0}
\fxwface{2}{3.25}{1}{0.5}{0}
\fxwface{2}{3.75}{1}{0.5}{0}
\fxwface{2}{4.125}{1}{0.25}{0}
\fxwface{2}{4.375}{1}{0.25}{0}
\fxwface{2}{4.625}{1}{0.25}{0}
\fxwface{2}{4.875}{1}{0.25}{0}
\fxwface{2}{4.125}{1.25}{0.25}{0}
\fxwface{2}{4.375}{1.25}{0.25}{0}
\fxwface{2}{4.625}{1.25}{0.25}{0}
\fxwface{2}{4.875}{1.25}{0.25}{0}
\fxwrow{0}{3.75}{0.25}{1}
\fxwrow{0}{4.125}{0.125}{-4}
\fxwrow{1}{4}{0}{8}
\fxwrow{1}{4}{0.25}{-4}
\fxwhanging{4}{0.25}
\fxwstencil{1}{4}{0.125}{-4}
\fxwstencil{1}{4}{0.375}{4}
\fxwstencil{2}{4.125}{0.25}{0.477}
\fxwstencil{2}{4.375}{0.25}{0.243}
\fxwstencil{2}{4.625}{0.25}{-0.0063}
\fxwstencil{2}{3.75}{0.5}{-1.08}
\fxwstencil{2}{4.125}{0.5}{0.235}
\fxwstencil{2}{4.375}{0.5}{0.143}
\fxwstencil{2}{4.625}{0.5}{-0.0107}
\fxwstencil{2}{4.125}{0.75}{0.144}
\fxwstencil{2}{4.375}{0.75}{0.094}
%
    \end{scope}
    \fill[tr wall] (3.0,-0.1) rectangle (5.0,0.0);
    \draw[black,line width=0.8pt] (3.0,0.0) -- (5.0,0.0);
    \draw[black!70,dash pattern=on 2pt off 1.5pt,line width=0.5pt] (4,-0.12) -- (4,1.3);
    \fxwnone
    \def\fxwhanging##1##2{\node[tr vertex,minimum size=3.6pt] at (##1,##2) {};}%
    % Generated by docs/techref/figures/generate.jl; do not edit.
\fxwcell{3}{0}{3.5}{0.5}{0}
\fxwcell{3.5}{0}{4}{0.5}{0}
\fxwcell{3}{0.5}{3.5}{1}{0}
\fxwcell{3.5}{0.5}{4}{1}{0}
\fxwcell{3}{1}{3.5}{1.5}{0}
\fxwcell{3.5}{1}{4}{1.5}{0}
\fxwcell{4}{0}{4.25}{0.25}{1}
\fxwcell{4.25}{0}{4.5}{0.25}{1}
\fxwcell{4.5}{0}{4.75}{0.25}{1}
\fxwcell{4.75}{0}{5}{0.25}{1}
\fxwcell{4}{0.25}{4.25}{0.5}{1}
\fxwcell{4.25}{0.25}{4.5}{0.5}{1}
\fxwcell{4.5}{0.25}{4.75}{0.5}{1}
\fxwcell{4.75}{0.25}{5}{0.5}{1}
\fxwcell{4}{0.5}{4.25}{0.75}{1}
\fxwcell{4.25}{0.5}{4.5}{0.75}{1}
\fxwcell{4.5}{0.5}{4.75}{0.75}{1}
\fxwcell{4.75}{0.5}{5}{0.75}{1}
\fxwcell{4}{0.75}{4.25}{1}{1}
\fxwcell{4.25}{0.75}{4.5}{1}{1}
\fxwcell{4.5}{0.75}{4.75}{1}{1}
\fxwcell{4.75}{0.75}{5}{1}{1}
\fxwcell{4}{1}{4.25}{1.25}{1}
\fxwcell{4.25}{1}{4.5}{1.25}{1}
\fxwcell{4.5}{1}{4.75}{1.25}{1}
\fxwcell{4.75}{1}{5}{1.25}{1}
\fxwface{1}{3}{0.25}{0.5}{0}
\fxwface{1}{3}{0.75}{0.5}{0}
\fxwface{1}{3}{1.25}{0.5}{0}
\fxwface{1}{3.5}{0.25}{0.5}{0}
\fxwface{1}{3.5}{0.75}{0.5}{0}
\fxwface{1}{3.5}{1.25}{0.5}{0}
\fxwface{1}{4}{0.125}{0.25}{1}
\fxwface{1}{4}{0.375}{0.25}{0}
\fxwface{1}{4}{0.625}{0.25}{0}
\fxwface{1}{4}{0.875}{0.25}{0}
\fxwface{1}{4}{1.125}{0.25}{0}
\fxwface{1}{4.25}{0.125}{0.25}{0}
\fxwface{1}{4.25}{0.375}{0.25}{0}
\fxwface{1}{4.25}{0.625}{0.25}{0}
\fxwface{1}{4.25}{0.875}{0.25}{0}
\fxwface{1}{4.25}{1.125}{0.25}{0}
\fxwface{1}{4.5}{0.125}{0.25}{0}
\fxwface{1}{4.5}{0.375}{0.25}{0}
\fxwface{1}{4.5}{0.625}{0.25}{0}
\fxwface{1}{4.5}{0.875}{0.25}{0}
\fxwface{1}{4.5}{1.125}{0.25}{0}
\fxwface{1}{4.75}{0.125}{0.25}{0}
\fxwface{1}{4.75}{0.375}{0.25}{0}
\fxwface{1}{4.75}{0.625}{0.25}{0}
\fxwface{1}{4.75}{0.875}{0.25}{0}
\fxwface{1}{4.75}{1.125}{0.25}{0}
\fxwface{1}{5}{0.125}{0.25}{0}
\fxwface{1}{5}{0.375}{0.25}{0}
\fxwface{1}{5}{0.625}{0.25}{0}
\fxwface{1}{5}{0.875}{0.25}{0}
\fxwface{1}{5}{1.125}{0.25}{0}
\fxwface{2}{4.125}{0.25}{0.25}{0}
\fxwface{2}{4.375}{0.25}{0.25}{0}
\fxwface{2}{4.625}{0.25}{0.25}{0}
\fxwface{2}{4.875}{0.25}{0.25}{0}
\fxwface{2}{3.25}{0.5}{0.5}{0}
\fxwface{2}{3.75}{0.5}{0.5}{0}
\fxwface{2}{4.125}{0.5}{0.25}{0}
\fxwface{2}{4.375}{0.5}{0.25}{0}
\fxwface{2}{4.625}{0.5}{0.25}{0}
\fxwface{2}{4.875}{0.5}{0.25}{0}
\fxwface{2}{4.125}{0.75}{0.25}{0}
\fxwface{2}{4.375}{0.75}{0.25}{0}
\fxwface{2}{4.625}{0.75}{0.25}{0}
\fxwface{2}{4.875}{0.75}{0.25}{0}
\fxwface{2}{3.25}{1}{0.5}{0}
\fxwface{2}{3.75}{1}{0.5}{0}
\fxwface{2}{4.125}{1}{0.25}{0}
\fxwface{2}{4.375}{1}{0.25}{0}
\fxwface{2}{4.625}{1}{0.25}{0}
\fxwface{2}{4.875}{1}{0.25}{0}
\fxwface{2}{4.125}{1.25}{0.25}{0}
\fxwface{2}{4.375}{1.25}{0.25}{0}
\fxwface{2}{4.625}{1.25}{0.25}{0}
\fxwface{2}{4.875}{1.25}{0.25}{0}
\fxwrow{0}{3.75}{0.25}{1}
\fxwrow{0}{4.125}{0.125}{-4}
\fxwrow{1}{4}{0}{8}
\fxwrow{1}{4}{0.25}{-4}
\fxwhanging{4}{0.25}
\fxwstencil{1}{4}{0.125}{-4}
\fxwstencil{1}{4}{0.375}{4}
\fxwstencil{2}{4.125}{0.25}{0.477}
\fxwstencil{2}{4.375}{0.25}{0.243}
\fxwstencil{2}{4.625}{0.25}{-0.0063}
\fxwstencil{2}{3.75}{0.5}{-1.08}
\fxwstencil{2}{4.125}{0.5}{0.235}
\fxwstencil{2}{4.375}{0.5}{0.143}
\fxwstencil{2}{4.625}{0.5}{-0.0107}
\fxwstencil{2}{4.125}{0.75}{0.144}
\fxwstencil{2}{4.375}{0.75}{0.094}
%
  }%
  \begin{subfigure}[b]{0.495\textwidth}
    \centering
    \begin{tikzpicture}[x=3.05cm,y=3.05cm]
      \fxwbackground{1}
      \fxwnone
      \def\fxwrow#1#2#3#4{%
        \ifnum#1=0
          \node[rectangle,draw=black,fill=trgreen!55,inner sep=0pt,minimum size=4.4pt] at (#2,#3) {};
          \node[tr value,anchor=south west] at (#2+0.015,#3+0.015) {$#4$};
        \else
          \node[diamond,draw=black,fill=trpurple!55,inner sep=0pt,minimum size=5.6pt] at (#2,#3) {};
          \node[tr value,anchor=west] at (#2+0.03,#3+0.05) {$#4$};
        \fi}%
      % Generated by docs/techref/figures/generate.jl; do not edit.
\fxwcell{3}{0}{3.5}{0.5}{0}
\fxwcell{3.5}{0}{4}{0.5}{0}
\fxwcell{3}{0.5}{3.5}{1}{0}
\fxwcell{3.5}{0.5}{4}{1}{0}
\fxwcell{3}{1}{3.5}{1.5}{0}
\fxwcell{3.5}{1}{4}{1.5}{0}
\fxwcell{4}{0}{4.25}{0.25}{1}
\fxwcell{4.25}{0}{4.5}{0.25}{1}
\fxwcell{4.5}{0}{4.75}{0.25}{1}
\fxwcell{4.75}{0}{5}{0.25}{1}
\fxwcell{4}{0.25}{4.25}{0.5}{1}
\fxwcell{4.25}{0.25}{4.5}{0.5}{1}
\fxwcell{4.5}{0.25}{4.75}{0.5}{1}
\fxwcell{4.75}{0.25}{5}{0.5}{1}
\fxwcell{4}{0.5}{4.25}{0.75}{1}
\fxwcell{4.25}{0.5}{4.5}{0.75}{1}
\fxwcell{4.5}{0.5}{4.75}{0.75}{1}
\fxwcell{4.75}{0.5}{5}{0.75}{1}
\fxwcell{4}{0.75}{4.25}{1}{1}
\fxwcell{4.25}{0.75}{4.5}{1}{1}
\fxwcell{4.5}{0.75}{4.75}{1}{1}
\fxwcell{4.75}{0.75}{5}{1}{1}
\fxwcell{4}{1}{4.25}{1.25}{1}
\fxwcell{4.25}{1}{4.5}{1.25}{1}
\fxwcell{4.5}{1}{4.75}{1.25}{1}
\fxwcell{4.75}{1}{5}{1.25}{1}
\fxwface{1}{3}{0.25}{0.5}{0}
\fxwface{1}{3}{0.75}{0.5}{0}
\fxwface{1}{3}{1.25}{0.5}{0}
\fxwface{1}{3.5}{0.25}{0.5}{0}
\fxwface{1}{3.5}{0.75}{0.5}{0}
\fxwface{1}{3.5}{1.25}{0.5}{0}
\fxwface{1}{4}{0.125}{0.25}{1}
\fxwface{1}{4}{0.375}{0.25}{0}
\fxwface{1}{4}{0.625}{0.25}{0}
\fxwface{1}{4}{0.875}{0.25}{0}
\fxwface{1}{4}{1.125}{0.25}{0}
\fxwface{1}{4.25}{0.125}{0.25}{0}
\fxwface{1}{4.25}{0.375}{0.25}{0}
\fxwface{1}{4.25}{0.625}{0.25}{0}
\fxwface{1}{4.25}{0.875}{0.25}{0}
\fxwface{1}{4.25}{1.125}{0.25}{0}
\fxwface{1}{4.5}{0.125}{0.25}{0}
\fxwface{1}{4.5}{0.375}{0.25}{0}
\fxwface{1}{4.5}{0.625}{0.25}{0}
\fxwface{1}{4.5}{0.875}{0.25}{0}
\fxwface{1}{4.5}{1.125}{0.25}{0}
\fxwface{1}{4.75}{0.125}{0.25}{0}
\fxwface{1}{4.75}{0.375}{0.25}{0}
\fxwface{1}{4.75}{0.625}{0.25}{0}
\fxwface{1}{4.75}{0.875}{0.25}{0}
\fxwface{1}{4.75}{1.125}{0.25}{0}
\fxwface{1}{5}{0.125}{0.25}{0}
\fxwface{1}{5}{0.375}{0.25}{0}
\fxwface{1}{5}{0.625}{0.25}{0}
\fxwface{1}{5}{0.875}{0.25}{0}
\fxwface{1}{5}{1.125}{0.25}{0}
\fxwface{2}{4.125}{0.25}{0.25}{0}
\fxwface{2}{4.375}{0.25}{0.25}{0}
\fxwface{2}{4.625}{0.25}{0.25}{0}
\fxwface{2}{4.875}{0.25}{0.25}{0}
\fxwface{2}{3.25}{0.5}{0.5}{0}
\fxwface{2}{3.75}{0.5}{0.5}{0}
\fxwface{2}{4.125}{0.5}{0.25}{0}
\fxwface{2}{4.375}{0.5}{0.25}{0}
\fxwface{2}{4.625}{0.5}{0.25}{0}
\fxwface{2}{4.875}{0.5}{0.25}{0}
\fxwface{2}{4.125}{0.75}{0.25}{0}
\fxwface{2}{4.375}{0.75}{0.25}{0}
\fxwface{2}{4.625}{0.75}{0.25}{0}
\fxwface{2}{4.875}{0.75}{0.25}{0}
\fxwface{2}{3.25}{1}{0.5}{0}
\fxwface{2}{3.75}{1}{0.5}{0}
\fxwface{2}{4.125}{1}{0.25}{0}
\fxwface{2}{4.375}{1}{0.25}{0}
\fxwface{2}{4.625}{1}{0.25}{0}
\fxwface{2}{4.875}{1}{0.25}{0}
\fxwface{2}{4.125}{1.25}{0.25}{0}
\fxwface{2}{4.375}{1.25}{0.25}{0}
\fxwface{2}{4.625}{1.25}{0.25}{0}
\fxwface{2}{4.875}{1.25}{0.25}{0}
\fxwrow{0}{3.75}{0.25}{1}
\fxwrow{0}{4.125}{0.125}{-4}
\fxwrow{1}{4}{0}{8}
\fxwrow{1}{4}{0.25}{-4}
\fxwhanging{4}{0.25}
\fxwstencil{1}{4}{0.125}{-4}
\fxwstencil{1}{4}{0.375}{4}
\fxwstencil{2}{4.125}{0.25}{0.477}
\fxwstencil{2}{4.375}{0.25}{0.243}
\fxwstencil{2}{4.625}{0.25}{-0.0063}
\fxwstencil{2}{3.75}{0.5}{-1.08}
\fxwstencil{2}{4.125}{0.5}{0.235}
\fxwstencil{2}{4.375}{0.5}{0.143}
\fxwstencil{2}{4.625}{0.5}{-0.0107}
\fxwstencil{2}{4.125}{0.75}{0.144}
\fxwstencil{2}{4.375}{0.75}{0.094}

      \node[tr small,anchor=north east,fill=white,fill opacity=0.85,text opacity=1] at (3.98,1.23) {$c$};
      \node[tr small,anchor=north west,fill=white,fill opacity=0.85,text opacity=1] at (4.27,0.23) {$c'$};
    \end{tikzpicture}
    \caption{Rows of $\strainmat$ that contain $\velo_f$}
    \label{fig:viscosity:stencil:rows}
  \end{subfigure}\hfill
  \begin{subfigure}[b]{0.495\textwidth}
    \centering
    \begin{tikzpicture}[x=3.05cm,y=3.05cm]
      \fxwbackground{2}
      \fxwnone
      \def\fxwstencil#1#2#3#4{%
        \ifnum#1=1 \node[circle,fill=trred,draw=black,line width=0.3pt,inner sep=0pt,minimum size=3.6pt] at (#2,#3) {};
        \else \node[circle,fill=trblue,draw=black,line width=0.3pt,inner sep=0pt,minimum size=3.6pt] at (#2,#3) {};\fi
        \node[tr value,anchor=south] at (#2,#3+0.03) {$#4$};}%
      % Generated by docs/techref/figures/generate.jl; do not edit.
\fxwcell{3}{0}{3.5}{0.5}{0}
\fxwcell{3.5}{0}{4}{0.5}{0}
\fxwcell{3}{0.5}{3.5}{1}{0}
\fxwcell{3.5}{0.5}{4}{1}{0}
\fxwcell{3}{1}{3.5}{1.5}{0}
\fxwcell{3.5}{1}{4}{1.5}{0}
\fxwcell{4}{0}{4.25}{0.25}{1}
\fxwcell{4.25}{0}{4.5}{0.25}{1}
\fxwcell{4.5}{0}{4.75}{0.25}{1}
\fxwcell{4.75}{0}{5}{0.25}{1}
\fxwcell{4}{0.25}{4.25}{0.5}{1}
\fxwcell{4.25}{0.25}{4.5}{0.5}{1}
\fxwcell{4.5}{0.25}{4.75}{0.5}{1}
\fxwcell{4.75}{0.25}{5}{0.5}{1}
\fxwcell{4}{0.5}{4.25}{0.75}{1}
\fxwcell{4.25}{0.5}{4.5}{0.75}{1}
\fxwcell{4.5}{0.5}{4.75}{0.75}{1}
\fxwcell{4.75}{0.5}{5}{0.75}{1}
\fxwcell{4}{0.75}{4.25}{1}{1}
\fxwcell{4.25}{0.75}{4.5}{1}{1}
\fxwcell{4.5}{0.75}{4.75}{1}{1}
\fxwcell{4.75}{0.75}{5}{1}{1}
\fxwcell{4}{1}{4.25}{1.25}{1}
\fxwcell{4.25}{1}{4.5}{1.25}{1}
\fxwcell{4.5}{1}{4.75}{1.25}{1}
\fxwcell{4.75}{1}{5}{1.25}{1}
\fxwface{1}{3}{0.25}{0.5}{0}
\fxwface{1}{3}{0.75}{0.5}{0}
\fxwface{1}{3}{1.25}{0.5}{0}
\fxwface{1}{3.5}{0.25}{0.5}{0}
\fxwface{1}{3.5}{0.75}{0.5}{0}
\fxwface{1}{3.5}{1.25}{0.5}{0}
\fxwface{1}{4}{0.125}{0.25}{1}
\fxwface{1}{4}{0.375}{0.25}{0}
\fxwface{1}{4}{0.625}{0.25}{0}
\fxwface{1}{4}{0.875}{0.25}{0}
\fxwface{1}{4}{1.125}{0.25}{0}
\fxwface{1}{4.25}{0.125}{0.25}{0}
\fxwface{1}{4.25}{0.375}{0.25}{0}
\fxwface{1}{4.25}{0.625}{0.25}{0}
\fxwface{1}{4.25}{0.875}{0.25}{0}
\fxwface{1}{4.25}{1.125}{0.25}{0}
\fxwface{1}{4.5}{0.125}{0.25}{0}
\fxwface{1}{4.5}{0.375}{0.25}{0}
\fxwface{1}{4.5}{0.625}{0.25}{0}
\fxwface{1}{4.5}{0.875}{0.25}{0}
\fxwface{1}{4.5}{1.125}{0.25}{0}
\fxwface{1}{4.75}{0.125}{0.25}{0}
\fxwface{1}{4.75}{0.375}{0.25}{0}
\fxwface{1}{4.75}{0.625}{0.25}{0}
\fxwface{1}{4.75}{0.875}{0.25}{0}
\fxwface{1}{4.75}{1.125}{0.25}{0}
\fxwface{1}{5}{0.125}{0.25}{0}
\fxwface{1}{5}{0.375}{0.25}{0}
\fxwface{1}{5}{0.625}{0.25}{0}
\fxwface{1}{5}{0.875}{0.25}{0}
\fxwface{1}{5}{1.125}{0.25}{0}
\fxwface{2}{4.125}{0.25}{0.25}{0}
\fxwface{2}{4.375}{0.25}{0.25}{0}
\fxwface{2}{4.625}{0.25}{0.25}{0}
\fxwface{2}{4.875}{0.25}{0.25}{0}
\fxwface{2}{3.25}{0.5}{0.5}{0}
\fxwface{2}{3.75}{0.5}{0.5}{0}
\fxwface{2}{4.125}{0.5}{0.25}{0}
\fxwface{2}{4.375}{0.5}{0.25}{0}
\fxwface{2}{4.625}{0.5}{0.25}{0}
\fxwface{2}{4.875}{0.5}{0.25}{0}
\fxwface{2}{4.125}{0.75}{0.25}{0}
\fxwface{2}{4.375}{0.75}{0.25}{0}
\fxwface{2}{4.625}{0.75}{0.25}{0}
\fxwface{2}{4.875}{0.75}{0.25}{0}
\fxwface{2}{3.25}{1}{0.5}{0}
\fxwface{2}{3.75}{1}{0.5}{0}
\fxwface{2}{4.125}{1}{0.25}{0}
\fxwface{2}{4.375}{1}{0.25}{0}
\fxwface{2}{4.625}{1}{0.25}{0}
\fxwface{2}{4.875}{1}{0.25}{0}
\fxwface{2}{4.125}{1.25}{0.25}{0}
\fxwface{2}{4.375}{1.25}{0.25}{0}
\fxwface{2}{4.625}{1.25}{0.25}{0}
\fxwface{2}{4.875}{1.25}{0.25}{0}
\fxwrow{0}{3.75}{0.25}{1}
\fxwrow{0}{4.125}{0.125}{-4}
\fxwrow{1}{4}{0}{8}
\fxwrow{1}{4}{0.25}{-4}
\fxwhanging{4}{0.25}
\fxwstencil{1}{4}{0.125}{-4}
\fxwstencil{1}{4}{0.375}{4}
\fxwstencil{2}{4.125}{0.25}{0.477}
\fxwstencil{2}{4.375}{0.25}{0.243}
\fxwstencil{2}{4.625}{0.25}{-0.0063}
\fxwstencil{2}{3.75}{0.5}{-1.08}
\fxwstencil{2}{4.125}{0.5}{0.235}
\fxwstencil{2}{4.375}{0.5}{0.143}
\fxwstencil{2}{4.625}{0.5}{-0.0107}
\fxwstencil{2}{4.125}{0.75}{0.144}
\fxwstencil{2}{4.375}{0.75}{0.094}

    \end{tikzpicture}
    \caption{Shear row of the hanging vertex}
    \label{fig:viscosity:stencil:shear}
  \end{subfigure}
  \caption{The strain-rate rows~\cref{eq:viscosity:rows} next to a level
    transition (dashed) and a no-slip wall (hatched), read from the code's
    operator on the walled composite fixture (coarse cells of width $\half$
    on the left). (a)~The \chkViscFigNormal\ normal rows (green squares, at
    the centres of the coarse cell $c$ and the fine cell $c'$) and the
    \chkViscFigShear\ shear rows (purple diamonds, at the vertices) whose
    entry for the fine segment $f$ (thick red) is not zero, with that entry:
    $\pm\aperture{f}\meas{f}/\meas{\opencell{c}}$ in the normal rows, the
    two-point difference to the wall and across the hanging vertex in the
    shear rows. (b)~The complete shear row of the hanging vertex (open circle):
    \chkViscFigStencil\ unknowns, $\velo_1$ (red) by the two-point difference
    along the fine segments, $\velo_2$ (blue) by the least-squares stencil
    of~\cref{eq:viscosity:shear}, since the coarse cell has no $\velo_2$
    unknown on the line through the vertex.}
  \label{fig:viscosity:stencil}
\end{figure}
